\documentclass[11pt]{article}
\usepackage[margin=2.7cm]{geometry}
\usepackage{amsmath,amssymb,booktabs,url}
\setlength{\parindent}{0pt}
\setlength{\parskip}{0.6\baselineskip}
\sloppy

\newcommand{\C}{\mathcal{C}}

\begin{document}

\begin{center}
{\large\bfseries Multiscale validity tests explored for \texttt{f1\_four\_methods.pdf}}
\end{center}

Figure \texttt{f1\_four\_methods.pdf} compares eleven change-point
detection methods on independent Gaussian data, all calibrated to a $1\%$
null false-positive rate at $n=1000$ observations (PELT and the ten SVP
variants all target the same rate, so their detection power is directly
comparable). Notation throughout follows the SVP framework of
\cite{runge2026svp} (Chapter~4, ``Smallest Valid Partitioning''): a
segment $y_{a..b}$ denotes the consecutive points $y_{a+1},\dots,y_b$, and
for $a<b$,
\[
\C^{\text{Gauss}}(y_{a..b}) \;=\; \sum_{i=a+1}^{b}\bigl(y_i-\overline{y}_{a..b}\bigr)^2,
\qquad \overline{y}_{a..b}=\frac{1}{b-a}\sum_{i=a+1}^{b}y_i,
\]
is the Gaussian segment cost (residual sum of squares, on the scale with
the customary $1/2$ factor removed, matching \cite{runge2026svp}'s own
convention). For $s<\tau<t$,
\[
f_{\mathrm{Gauss}}(y_{s..t})
= 2f_{\mathrm{LR}}(y_{s..t})
= \max_{s<\tau<t}\Bigl\{\C^{\text{Gauss}}(y_{s..t})-\C^{\text{Gauss}}(y_{s..\tau})-\C^{\text{Gauss}}(y_{\tau..t})\Bigr\}
\]
is the single-anchored Gaussian generalized-likelihood-ratio statistic
(start $s$ and end $t$ fixed, maximized only over the internal split
$\tau$): twice the log-likelihood-ratio for one versus two segments on
$(s,t]$, and the Gaussian specialization of the cost-based validity test
$f_{OP}(y_{a..b})=\max_{a<u<b}\{\C(y_{a..b})-\C(y_{a..u})-\C(y_{u..b})\}$
of \cite{runge2026svp} (\S``Link to the OP method''). Following the same
construction, write $\Delta(a,u,b)$ -- a local shorthand for this
document, not a symbol used in \cite{runge2026svp} itself -- for the
bracket inside that maximum at one specific, unmaximized split $a<u<b$,
\[
\Delta(a,u,b) \;=\; \C^{\text{Gauss}}(y_{a..b})-\C^{\text{Gauss}}(y_{a..u})-\C^{\text{Gauss}}(y_{u..b}),
\]
so that $f_{\mathrm{Gauss}}(y_{s..t})=\max_{s<u<t}\Delta(s,u,t)$. It is
also useful to recall the multiscale aggregates of \cite{runge2026svp}
built from any single-change test $f$,
\[
F_{\mathrm{right}}(y_{s..t})=\max_{s<\tau\le t}f(y_{s..\tau}),\quad
F_{\mathrm{left}}(y_{s..t})=\max_{s\le \tau<t}f(y_{\tau..t}),\quad
F_{\mathrm{both}}(y_{s..t})=\max_{s\le \tau_1<\tau_2\le t}f(y_{\tau_1..\tau_2}),
\]
which are, respectively, $\gamma$-stable, $\gamma^{-1}$-stable, and both
-- the property that makes discard-on-failure (\texttt{"right"}) and
inclusive left-pruning (\texttt{"both"}) pruning exact \emph{when the
validity test itself is} $F_{\mathrm{right}}$ or $F_{\mathrm{both}}$.
$n$ denotes the total series length and $\ell=t-s$ a candidate segment's
own length. All ten SVP variants search for the partition with the
\emph{fewest} segments passing their validity test, breaking ties by
minimal total within-segment sum of squared errors.

\section*{1. The six methods in \texttt{f1\_six\_methods.pdf}}

The three fixed baselines (PELT, SVP both, SVP SMUCE) plus the three
best-performing \texttt{"both"}-family variants, selected by mean F1
across all scenarios and jump sizes (and, separately, restricted to the
harder jump $\le1.0$ regime -- both rankings agree): \texttt{both\_rmin},
\texttt{both\_correction}, and \texttt{both\_mosum}. These six are
collected on their own in \texttt{f1\_six\_methods.pdf}; the remaining
five variants (\S2) are shown only in the full \texttt{f1\_four\_methods.pdf}.

\begin{itemize}

\item \textbf{PELT.} Exact minimizer of a penalized Gaussian cost,
\[
\min_{m,\,\tau_1<\dots<\tau_m}\ \sum_{i=0}^{m}\C^{\text{Gauss}}(y_{\tau_i..\tau_{i+1}})
\;+\;\beta m,
\qquad \tau_0=0,\ \tau_{m+1}=n,\ \beta=c\log n,
\]
solved exactly by the pruned dynamic program of \cite{killick2012}, and
matching \cite{runge2026svp}'s own $Q_n(\tau;y)+\gamma K_n(\tau)$
objective (no $1/2$ factor, as above); $f_{\mathrm{Gauss}}(y_{s..t})$ is
exactly the drop in this same cost from optimally splitting $(s,t]$ once.
\emph{Calibration:} $c$ found by simulation so PELT itself reaches a $1\%$
null FPR at $n=1000$ (not the usual BIC-style $c=1$ or $2$), for a fair
comparison against the SVP variants. \emph{Reference:}
\cite{killick2012}.

\item \textbf{SVP both.} Smallest partition with every segment
$(s,t]$ satisfying
\[
f_{\mathrm{Gauss}}(y_{s..t}) \;\le\; c\log n.
\]
\emph{Calibration:} $c$ found by simulation to a $1\%$ null FPR under
\texttt{"both"} pruning (discard a candidate on failure, plus inclusive
left-pruning of older candidates sharing the same endpoint) -- applied as
a heuristic shortcut to the raw, single-anchored $f_{\mathrm{Gauss}}$, not
to $F_{\mathrm{both}}$ itself.
\emph{Reference:} the FOCUS statistic and the
\texttt{"both"}/\texttt{"right"}/\texttt{"none"} functional-pruning rules
are from \cite{romano2023}; the SVP smallest-valid-partition search and
the $F_{\mathrm{right}}/F_{\mathrm{left}}/F_{\mathrm{both}}$ framework are
from \cite{runge2026svp}.

\item \textbf{SVP SMUCE.} Smallest partition with every segment $(s,t]$
of length $\ell$ satisfying: there exists $\theta$ such that every
sub-window $(s',t']\subseteq(s,t]$ of length $\ell'$ obeys
\[
\bigl|\overline{y}_{s'..t'}-\theta\bigr| \;\le\;
\frac{1}{\sqrt{\ell'}}\Bigl(q+\sqrt{2\bigl(1+\log(n/\ell')\bigr)}\Bigr).
\]
\emph{Calibration:} $q$ found by simulation to a $1\%$ null FPR.
\emph{Reference:} the multiscale constraint (SMUCE) is from
\cite{frick2014}; \cite{runge2026svp} notes $F_{\mathrm{both}}$, evaluated
with a fixed-scale validity test $f$, is precisely the multiscale
statistic underlying SMUCE.

\item \textbf{SVP both\_rmin.} $f_{\mathrm{Gauss}}$'s inner maximum
restricted to splits whose shorter arm exceeds a minimum length,
\[
\max_{\tau:\,r_\tau\ge r_{\min}} \Delta(s,\tau,t) \;\le\; \gamma.
\]
\emph{Calibration:} grid search over $r_{\min}$, then $\gamma$ bisected to
a $1\%$ null FPR at each ($r_{\min}=30$, $\gamma=25.313$) -- again
dispatched via the \texttt{"both"}-pruning heuristic on this restricted
raw $f_{\mathrm{Gauss}}$, not on an $F_{\mathrm{both}}$ aggregate of it.
\emph{Reference:} none found; an internal \texttt{svpChange}
design choice.

\item \textbf{SVP both\_correction.} Every sub-interval $[a,b]\subseteq[s,t]$
tested, at every internal split $u$, against a length-corrected penalty --
the same double-max family as $F_{\mathrm{both}}$, but with the penalty
applied \emph{before} maximizing over $u$, so the unmaximized $\Delta$ is
needed rather than $f_{\mathrm{Gauss}}$ itself:
\[
\max_{s\le a<u<b\le t}\Bigl\{\sqrt{\Delta(a,u,b)} - p_n(u-a,\,b-u)\Bigr\} \;\le\; q,
\qquad p_n(l,r)=\sqrt{2\log\!\bigl(n(l+r)/(lr)\bigr)}.
\]
\emph{Calibration:} $q=2.623$ found by simulation to a $1\%$ null FPR.
\emph{Reference:} \texttt{SVP\_consistency\_sharp.tex} (unpublished draft);
the penalty $p_n$ is derived from the simultaneous multiscale CUSUM bound
of \cite{verzelen2020}.

\item \textbf{SVP both\_mosum.} The single-anchored statistic evaluated
at a single, fixed split -- the candidate's own midpoint -- instead of
maximized over every internal split as $f_{\mathrm{Gauss}}$ is: for a
candidate of current length $\ell$, with $\tau_{\mathrm{mid}}=s+\lfloor
\ell/2\rfloor$,
\[
\Delta(s,\tau_{\mathrm{mid}},t) \;\le\; \gamma,
\qquad \text{segments with } \ell<20 \text{ always valid}.
\]
Unlike every other method here, no maximum over $\tau$ is taken at all --
a single two-sample comparison per candidate, at $O(1)$ cost instead of
the $O(\ell)$ or worse scans used elsewhere in this document -- and
segments shorter than $20$ are declared valid unconditionally (a hard
floor, not an asymptotic one: below that length a single fixed midpoint
split leaves too few points on one side to be informative).
\emph{No multiscale penalty:} unlike every other method in this document,
$\gamma$ does not depend on $\ell$. This is not an oversight: because no
maximum over $\tau$ is taken, $\Delta(s,\tau_{\mathrm{mid}},t)$ is
\emph{exactly} $\chi^2_1$-distributed at every length $\ell$ (verified by
simulation: mean $\approx1$ at $\ell=20,\dots,1000$), unlike
$f_{\mathrm{Gauss}}$, whose own maximum over $\sim\ell$ positions makes its
null distribution grow like $2\log\ell$ -- the very reason
\texttt{both\_calibrated}, \texttt{both\_lengthonly}, and
\texttt{both\_mosum\_sqrt} all need a length-adaptive threshold. A flat $\gamma$
is therefore the right \emph{shape} here. A large multiplicity still has
to be absorbed, though: testing the same candidate at every length it
survives under \texttt{"both"} pruning means a single test's marginal
exceedance probability at $\gamma=23.32$ ($\chi^2_1$'s tail,
$1.4\times10^{-6}$) is far below the realized $1\%$ aggregate FPR -- an
implied effective $K\approx7300$ tests, all absorbed by the one calibrated
constant rather than any parametric union-bound formula.
\emph{Calibration:} $\gamma=23.32$, a flat threshold found by directly
bisecting the empirical null rate to $1\%$ under \texttt{"both"} pruning
(true FPR $0.96\%$ at $8000$ replicates) -- the same heuristic shortcut on
this raw statistic as every other method here. \emph{Power:} F1 at jump
$0.6$ on \texttt{updown} is $0.109$, on par with the $f_{\mathrm{Gauss}}$-based
methods ($0.07$--$0.13$) despite testing only one split per candidate
instead of maximizing over all of them: because no maximum over $\tau$ is
taken, this statistic needs no correction for scanning multiple splits
within a candidate (unlike $f_{\mathrm{Gauss}}$, whose own internal
maximization already inflates its null distribution), so the threshold
this test needs is comparatively mild, recovering most of the power the
internal split-search would otherwise have bought. \emph{Reference:} none
found; an internal \texttt{svpChange} design choice.

\end{itemize}

\section*{2. Five further variants (also in \texttt{f1\_four\_methods.pdf})}

Shown only in the full figure, not in \texttt{f1\_six\_methods.pdf}: three
further ways of recalibrating or restructuring the methods above
(\texttt{both\_calibrated Bonferroni}, \texttt{both\_empirical},
\texttt{both\_lengthonly}), and two statistics that replace
$f_{\mathrm{Gauss}}$ itself (\texttt{both\_mosum\_sqrt},
\texttt{both\_mosum\_adaptive}).

\begin{itemize}

\item \textbf{SVP both\_calibrated Bonferroni.} $f_{\mathrm{Gauss}}$
thresholded at each candidate's own current length $\ell$,
\[
f_{\mathrm{Gauss}}(y_{s..t}) \;\le\; A\log\ell + B\log\log\ell + C.
\]
\emph{Calibration:} $(A,B)=(1.683,\,6.209)$ are the original fitted
values, held fixed; $C=15.5$ (up from the original fit's $4.939$) is found
by directly bisecting the empirical null rate to $1\%$ -- dispatched via
the same \texttt{"both"}-pruning heuristic on raw $f_{\mathrm{Gauss}}$ as
plain \texttt{"both"}, not on its $F_{\mathrm{both}}$ aggregate.
\emph{Reference:} motivated by the union-bound principle of the classical
Bonferroni correction \cite{bonferroni1936}; the specific value of $C$ is
found empirically, with no external reference for that step.

\item \textbf{SVP both\_empirical.} No parametric formula: a simulated,
time-varying threshold on $F_{\mathrm{both}}$, the TRUE aggregate
\texttt{"both"} statistic (start fixed at $0$),
\[
F_t \;>\; E_t + c\sqrt{V_t},
\qquad F_t = F_{\mathrm{both}}(y_{0..t}) = \max_{0\le a<b\le t} f_{\mathrm{Gauss}}(y_{a..b}),
\]
the running maximum, over every sub-segment of the first $t$ observations,
of $f_{\mathrm{Gauss}}$ -- the statistic \cite{runge2026svp} proves is
exactly $\gamma$- and $\gamma^{-1}$-stable, matching what \texttt{"both"}
pruning's exactness requires. At runtime each candidate still
compares its own raw $f_{\mathrm{Gauss}}(y_{s..t})$ against the shared
$\gamma(t)=E_t+c\sqrt{V_t}$ (the same heuristic shortcut as every other
\texttt{"both"}-pruned method here), but this is the one method whose
\emph{threshold} is calibrated directly against $F_{\mathrm{both}}$'s own
simulated distribution, rather than by bisecting the real dispatched
algorithm's FPR blind to that distinction. \emph{Calibration:} $E_t,V_t$
(per-length mean and variance of $F_{\mathrm{both}}(y_{0..t})$, obtained
by a running cummax of the per-length $F_{\mathrm{left}}$ trajectory,
since $F_{\mathrm{both}}(y_{0..t})=\max_{t'\le t}F_{\mathrm{left}}(y_{0..t'})$)
and $c$ (the joint running-maximum quantile) all estimated empirically
from $20{,}000$ null trajectories ($c=5.187$). \emph{Reference:}
$F_{\mathrm{both}}$ is from \cite{runge2026svp}; the empirical threshold
construction is original to this document.

\item \textbf{SVP both\_lengthonly.} $f_{\mathrm{Gauss}}$ thresholded by a
length-only multiscale correction, depending on the candidate's own
current length $\ell$ but not on the split point:
\[
f_{\mathrm{Gauss}}(y_{s..t}) \;\le\; \Bigl[q+\sqrt{2\log(en/\ell)}\Bigr]^2.
\]
\emph{Calibration:} $q=3.30$ found by directly bisecting the empirical
null rate to a $1\%$ FPR under \texttt{"both"} pruning -- the same
heuristic shortcut on raw $f_{\mathrm{Gauss}}$ as every other
\texttt{"both"}-pruned method here. Unlike the original proposal, every
length $\ell$ is scanned (no declared length subset) and a single shared
$q$ is used (no per-length optimization). \emph{Reference:} the
$\sqrt{2\log(en/\ell)}$ multiscale correction follows the principle of
\cite{dumbgen2001} and its change-point form, SMUCE, in \cite{frick2014};
the length-adaptive threshold itself is the
\texttt{SVP\_validity\_rules.tex} proposal (unpublished draft,
V.~Runge).

\item \textbf{SVP both\_mosum\_sqrt.} Replaces $f_{\mathrm{Gauss}}$ with a
moving-sum (MOSUM) statistic at bandwidth $G$: for a candidate of current
length $\ell$,
\[
T_k(G) \;=\; \frac{1}{\sqrt{2G\sigma^2}}\Biggl(\sum_{i=k+1}^{k+G}y_i-\sum_{i=k-G+1}^{k}y_i\Biggr),
\qquad G\le k\le \ell-G,
\]
compares the mean of the $G$ points to the right of $k$ against the mean
of the $G$ points to the left -- a fixed-width local two-sample test,
unlike $f_{\mathrm{Gauss}}$'s split-optimized comparison. The multiscale
penalty accounting for the candidate's own length $\ell$ is the exact
asymptotic boundary of \cite{eichinger2018}, Theorem~2.1: with
$G=G(\ell)=\max(2,\lfloor\sqrt\ell\rfloor)$ (satisfying their rate
condition $\ell/G\to\infty$) and $x=\ell/G$,
\[
a(x)=\sqrt{2\log x},\qquad
b(x)=2\log x+\tfrac12\log\log x+\log\tfrac32-\tfrac12\log\pi,
\]
\[
\max_{G\le k\le \ell-G}T_k(G)^2 \;\le\;
\left[\frac{b(x)+q_\alpha}{a(x)}\right]^2,
\qquad q_\alpha=-\log\!\Bigl(-\tfrac12\log(1-\alpha)\Bigr),
\]
the $(1-\alpha)$-quantile of the Gumbel limit
$a(n/G)\max_k|T_k(G)|-b(n/G)\xrightarrow{\mathcal D}\Gamma$,
$P(\Gamma\le x)=\exp(-2\exp(-x))$, with the total series length $n$
replaced by the candidate's own $\ell$ to make the threshold length-adaptive.
\emph{Calibration:} at $\gamma=0$ (the bare theoretical boundary, $\alpha=1\%$)
the empirical null FPR is already $1.4\%$, validating the asymptotic
derivation; a single additive correction $\gamma=1.625$, found by directly
bisecting the empirical null rate (same convention as
\texttt{both\_calibrated}: $\gamma$ shifts the squared statistic minus the
theoretical boundary), gives a true $1\%$ FPR ($0.76\%$ at $8000$
replicates) -- dispatched via the same \texttt{"both"}-pruning heuristic
on this raw statistic as every other method here. \emph{Power:}
substantially below every other method at small and moderate jump sizes
(F1 at jump $0.6$ on \texttt{updown} is $0.015$, versus $0.07$--$0.13$ for
the $f_{\mathrm{Gauss}}$-based methods), closing the gap only at large
jumps; this is a real property of the fixed-bandwidth MOSUM scan, not a
calibration artifact -- $T_k(G)$ compares two length-$G$ windows without
optimizing the split location the way $f_{\mathrm{Gauss}}$ does, so it is
mismatched whenever the true change does not fall near a well-covered $k$,
or $G$ is far from the true arm lengths. \emph{Reference:} the use of
MOSUM statistics for testing the i.i.d.\ multiple-change hypothesis is
from \cite{huskova2001}; the exact multiscale boundary used here is
\cite{eichinger2018}, Theorem~2.1.

\item \textbf{SVP both\_mosum\_adaptive.} The length-adaptive version
of \texttt{both\_mosum}: the raw statistic
$\Delta(s,\tau_{\mathrm{mid}},t)$ is unchanged (still a single, fixed
midpoint comparison), but the flat $\gamma$ is replaced by a
length-dependent threshold motivated as follows. Unlike
$f_{\mathrm{Gauss}}$'s own running maximum, which is non-stationary and
needs a single-log, Darling-Erd\H{o}s-type boundary
$\sqrt{2\log\ell}$ (the shape underlying
\texttt{both\_calibrated}/\texttt{both\_lengthonly}/\texttt{both\_mosum\_sqrt}),
the signed midpoint statistic $Z_{\mathrm{mid}}(\ell)$ is \emph{already
stationary} after a log-time change: its correlation across two lengths
$\ell<\ell'$ depends only on the ratio $r=\ell'/\ell$,
\[
\mathrm{Corr}\bigl(Z_{\mathrm{mid}}(\ell),Z_{\mathrm{mid}}(\ell')\bigr)
=\begin{cases}(2-r)/\sqrt{r}, & 1\le r<2,\\ 0, & r\ge2,\end{cases}
\]
a compactly supported correlation (derived from the Brownian-motion
limit, confirmed by simulation to within $0.01$). Controlling the running
maximum of such a stationary, short-range-dependent process under
\emph{repeated testing as the sample size grows} is the classical setting
of the law-of-the-iterated-logarithm literature on repeated significance
tests (\cite{darling1956}'s general boundary-crossing theorem;
\cite{darlingrobbins1967} on sequential/repeated significance testing),
which gives a much \emph{gentler}, double-log threshold growth than the
single-log corrections above. This motivates
\[
\Delta(s,\tau_{\mathrm{mid}},t) \;\le\;
\Bigl[q+\sqrt{2\log\bigl(\log(\ell/20)+e\bigr)}\Bigr]^2,
\]
(the ``$+e$'' keeps the inner argument $\ge1$ at $\ell=20$, avoiding
$\log(0)$). As with \texttt{both\_lengthonly} and \texttt{both\_mosum\_sqrt},
the exact extremal constant for this specific compactly-supported
correlation structure is not derived here in closed form; only the
general double-log \emph{shape} is taken from the theory, with $q$
calibrated empirically. \emph{Calibration:} $q=3.304$, found by directly
bisecting the empirical null rate to $1\%$ (true FPR $0.84\%$ at $8000$
replicates). \emph{Power:} F1 at jump $0.6$ on \texttt{updown} is
$0.102$, essentially unchanged from (very slightly below)
\texttt{both\_mosum}'s flat-threshold $0.109$: the length-adaptive
threshold is more liberal at short $\ell$ and more conservative at long
$\ell$ than the flat one ($22.3$ vs.\ $23.3$ at $\ell=20$; $27.6$ vs.\
$23.3$ at $\ell=1000$), and on \texttt{updown} this redistribution is
roughly power-neutral. \emph{Reference:} the general boundary-crossing
theorem is \cite{darling1956}; the repeated-significance-testing
framework motivating the double-log shape is
\cite{darlingrobbins1967}; the specific threshold formula and the
compactly-supported correlation derivation are original to this
document.

\end{itemize}

\begin{thebibliography}{9}

\bibitem{bonferroni1936}
C.~E.~Bonferroni.
\newblock Teoria statistica delle classi e calcolo delle probabilit\`a.
\newblock \emph{Pubblicazioni del R.~Istituto Superiore di Scienze
Economiche e Commerciali di Firenze}, 8:3--62, 1936.

\bibitem{darling1956}
D.~A.~Darling and P.~Erd\H{o}s.
\newblock A limit theorem for the maximum of normalized sums of
independent random variables.
\newblock \emph{Duke Mathematical Journal}, 23(1):143--155, 1956.

\bibitem{darlingrobbins1967}
D.~A.~Darling and H.~Robbins.
\newblock Iterated logarithm inequalities.
\newblock \emph{Proceedings of the National Academy of Sciences},
57(5):1188--1192, 1967.

\bibitem{dumbgen2001}
L.~D\"umbgen and V.~G.~Spokoiny.
\newblock Multiscale testing of qualitative hypotheses.
\newblock \emph{The Annals of Statistics}, 29(1):124--152, 2001.
\newblock \url{https://doi.org/10.1214/aos/996986504}.

\bibitem{eichinger2018}
B.~Eichinger and C.~Kirch.
\newblock A MOSUM procedure for the estimation of multiple random change
points.
\newblock \emph{Bernoulli}, 24(1):526--564, 2018.
\newblock \url{https://doi.org/10.3150/16-BEJ887}.

\bibitem{frick2014}
K.~Frick, A.~Munk and H.~Sieling.
\newblock Multiscale change point inference.
\newblock \emph{Journal of the Royal Statistical Society: Series B},
76(3):495--580, 2014.
\newblock \url{https://doi.org/10.1111/rssb.12047}.

\bibitem{huskova2001}
M.~Hu\v{s}kov\'a and A.~Slab\'y.
\newblock Permutation tests for multiple changes.
\newblock \emph{Kybernetika (Prague)}, 37(5):605--622, 2001.

\bibitem{killick2012}
R.~Killick, P.~Fearnhead and I.~A.~Eckley.
\newblock Optimal detection of changepoints with a linear computational cost.
\newblock \emph{Journal of the American Statistical Association},
107(500):1590--1598, 2012.
\newblock \url{https://doi.org/10.1080/01621459.2012.737745}.

\bibitem{romano2023}
G.~Romano, I.~A.~Eckley, P.~Fearnhead and G.~Rigaill.
\newblock Fast online changepoint detection via functional pruning CUSUM
statistics.
\newblock \emph{Journal of Machine Learning Research}, 24(81):1--36, 2023.
\newblock \url{https://www.jmlr.org/papers/v24/21-1230.html}.

\bibitem{runge2026svp}
V.~Runge, A.~Kostic, A.~Combeau and G.~Romano.
\newblock Exact multiple change-point detection via smallest valid
partitioning.
\newblock arXiv:2602.04322 [stat.ME], 2026.
\newblock \url{https://arxiv.org/abs/2602.04322}.

\bibitem{verzelen2020}
N.~Verzelen, M.~Fromont, M.~Lerasle and P.~Reynaud-Bouret.
\newblock Optimal change-point detection and localization.
\newblock \emph{Annals of Statistics}, 51(4):1586--1610, 2023.
\newblock \url{https://doi.org/10.1214/23-AOS2297}.
\newblock Detailed statements and proofs cited here use arXiv:2010.11470v2,
2020, \url{https://arxiv.org/abs/2010.11470v2}.

\end{thebibliography}

\end{document}
